%!TEX root = ../../main.tex
\section{교전 전 상태와 입력의 구성}
\label{sec:data-inputs}

어떤 질의 시각의 \emph{상태}는 그 시각 이전(포함)의 마지막 1분 프레임과 그 시각까지의 모든 사건으로 재구성한다(닫힌 포함, 보간 없음, 프레임 나이 저장). 승률 평가기 $\Vhat$의 상태 표현 Expanded\dimExpanded{}은 \dimNumeric{}개의 수치 특징과 \dimChampSlots{}개의 챔피언 ID 슬롯으로 이루어진다. 수치 블록은 프레임 기반의 자원·성장·체력 상태, 사건 이력, 지도 앵커, 시간 관련 열을 포함하며, 열 단위의 원천과 계보는 저장소의 특징 매니페스트에 기록되어 있다. 본 논문은 열을 하나씩 나열하지 않고 블록 수준에서 기술한다.

방향 예측기 $q$의 입력은 같은 \dimNumeric{}개의 수치 특징에 교전 전 승률 $\ppre = \Vhat(\xpre)$를 더한 \dimInputQ{}개이며, 챔피언 ID 슬롯은 제외한다. 진단용 LightGBM 후보는 같은 수치 블록에 챔피언 ID 10개를 범주형으로 더한 \dimInputLgbm{}개 입력을 썼으므로, 두 학습기는 동일한 입력에서 비교된 것이 아니다. 이 사실은 M-RQ3의 ``학습기 선택'' 절이 동결 설계로는 완전히 뒷받침되지 않음을 뜻하며, 제 \ref{sec:pred-learner} 절에서 다시 밝힌다.

v3.3 프리셋은 입력 특징에 관한 세 설정을 고정한다. 경기 시각은 45분으로 나누어 1에서 자르고(경기 자체의 길이로 정규화하지 않는다), 지도 앵커는 예측 시점 이전에 파괴된 구조물만 반영하며, 마지막 프레임의 나이를 추가 열로 둔다. 첫 번째와 두 번째 설정은 선행 연구의 파이프라인이 각 경기의 길이로 시간을 정규화하고 나중에 획득된 구조물·오브젝트를 앵커에 포함했던 누출을 고친 것이며, 이 누출이 사후 분석의 ``규모 기울기'' 결론을 만들어 냈다(제 \ref{sec:eng-scale} 절).
